\documentclass[12pt,a4paper]{article}
\usepackage[UTF8]{ctex}
\usepackage{amsmath,amssymb,amsthm}
\usepackage{geometry}
\usepackage{graphicx}
\usepackage{hyperref}
\usepackage{bm}

\geometry{margin=2.5cm}

\title{矩阵特征值与特征向量的计算}
\author{Modern Robotics}
\date{}

\begin{document}

\maketitle

\section{问题}

给定矩阵：
\begin{equation}
\bm{A} = \begin{bmatrix} 1 & 1 \\ -4 & 4 \end{bmatrix}
\end{equation}

求矩阵$\bm{A}$的特征值和特征向量。

\section{特征值计算}

特征值通过求解特征方程$\det(\bm{A} - \lambda \bm{I}) = 0$得到。

\subsection{步骤1：构造特征矩阵}

\begin{align}
\bm{A} - \lambda \bm{I} &= \begin{bmatrix} 1 & 1 \\ -4 & 4 \end{bmatrix} - \lambda \begin{bmatrix} 1 & 0 \\ 0 & 1 \end{bmatrix} \\
&= \begin{bmatrix} 1 - \lambda & 1 \\ -4 & 4 - \lambda \end{bmatrix}
\end{align}

\subsection{步骤2：计算特征多项式}

\begin{align}
\det(\bm{A} - \lambda \bm{I}) &= \det\begin{bmatrix} 1 - \lambda & 1 \\ -4 & 4 - \lambda \end{bmatrix} \\
&= (1 - \lambda)(4 - \lambda) - (1)(-4) \\
&= (1 - \lambda)(4 - \lambda) + 4 \\
&= 4 - 4\lambda - \lambda + \lambda^2 + 4 \\
&= \lambda^2 - 5\lambda + 8
\end{align}

\subsection{步骤3：求解特征方程}

令特征多项式等于零：
\begin{equation}
\lambda^2 - 5\lambda + 8 = 0
\label{eq:characteristic}
\end{equation}

使用二次公式求解：
\begin{align}
\lambda &= \frac{5 \pm \sqrt{25 - 32}}{2} \\
&= \frac{5 \pm \sqrt{-7}}{2} \\
&= \frac{5 \pm i\sqrt{7}}{2}
\end{align}

因此，特征值为：
\begin{align}
\lambda_1 &= \frac{5 + i\sqrt{7}}{2} \\
\lambda_2 &= \frac{5 - i\sqrt{7}}{2}
\end{align}

\section{特征向量计算}

对于每个特征值，求解齐次线性方程组$(\bm{A} - \lambda \bm{I})\bm{v} = \bm{0}$。

\subsection{特征向量$\bm{v}_1$对应特征值$\lambda_1$}

求解方程组：
\begin{equation}
(\bm{A} - \lambda_1 \bm{I})\bm{v}_1 = \bm{0}
\end{equation}

即：
\begin{align}
\begin{bmatrix} 1 - \lambda_1 & 1 \\ -4 & 4 - \lambda_1 \end{bmatrix} \begin{bmatrix} v_{11} \\ v_{12} \end{bmatrix} &= \begin{bmatrix} 0 \\ 0 \end{bmatrix}
\end{align}

代入$\lambda_1 = \frac{5 + i\sqrt{7}}{2}$：
\begin{align}
1 - \lambda_1 &= 1 - \frac{5 + i\sqrt{7}}{2} = \frac{2 - 5 - i\sqrt{7}}{2} = \frac{-3 - i\sqrt{7}}{2} \\
4 - \lambda_1 &= 4 - \frac{5 + i\sqrt{7}}{2} = \frac{8 - 5 - i\sqrt{7}}{2} = \frac{3 - i\sqrt{7}}{2}
\end{align}

因此方程组为：
\begin{equation}
\begin{bmatrix} \frac{-3 - i\sqrt{7}}{2} & 1 \\ -4 & \frac{3 - i\sqrt{7}}{2} \end{bmatrix} \begin{bmatrix} v_{11} \\ v_{12} \end{bmatrix} = \begin{bmatrix} 0 \\ 0 \end{bmatrix}
\end{equation}

从第一个方程：
\begin{align}
\frac{-3 - i\sqrt{7}}{2} v_{11} + v_{12} &= 0 \\
v_{12} &= \frac{3 + i\sqrt{7}}{2} v_{11}
\end{align}

取$v_{11} = 2$（方便计算），则：
\begin{align}
v_{12} &= \frac{3 + i\sqrt{7}}{2} \cdot 2 = 3 + i\sqrt{7}
\end{align}

因此，特征向量为：
\begin{equation}
\bm{v}_1 = \begin{bmatrix} 2 \\ 3 + i\sqrt{7} \end{bmatrix}
\end{equation}

或者，更一般的形式（任意非零标量倍数）：
\begin{equation}
\bm{v}_1 = c_1 \begin{bmatrix} 2 \\ 3 + i\sqrt{7} \end{bmatrix}, \quad c_1 \neq 0
\end{equation}

\subsection{特征向量$\bm{v}_2$对应特征值$\lambda_2$}

求解方程组：
\begin{equation}
(\bm{A} - \lambda_2 \bm{I})\bm{v}_2 = \bm{0}
\end{equation}

代入$\lambda_2 = \frac{5 - i\sqrt{7}}{2}$：
\begin{align}
1 - \lambda_2 &= 1 - \frac{5 - i\sqrt{7}}{2} = \frac{2 - 5 + i\sqrt{7}}{2} = \frac{-3 + i\sqrt{7}}{2} \\
4 - \lambda_2 &= 4 - \frac{5 - i\sqrt{7}}{2} = \frac{8 - 5 + i\sqrt{7}}{2} = \frac{3 + i\sqrt{7}}{2}
\end{align}

因此方程组为：
\begin{equation}
\begin{bmatrix} \frac{-3 + i\sqrt{7}}{2} & 1 \\ -4 & \frac{3 + i\sqrt{7}}{2} \end{bmatrix} \begin{bmatrix} v_{21} \\ v_{22} \end{bmatrix} = \begin{bmatrix} 0 \\ 0 \end{bmatrix}
\end{equation}

从第一个方程：
\begin{align}
\frac{-3 + i\sqrt{7}}{2} v_{21} + v_{22} &= 0 \\
v_{22} &= \frac{3 - i\sqrt{7}}{2} v_{21}
\end{align}

取$v_{21} = 2$，则：
\begin{align}
v_{22} &= \frac{3 - i\sqrt{7}}{2} \cdot 2 = 3 - i\sqrt{7}
\end{align}

因此，特征向量为：
\begin{equation}
\bm{v}_2 = \begin{bmatrix} 2 \\ 3 - i\sqrt{7} \end{bmatrix}
\end{equation}

或者，更一般的形式：
\begin{equation}
\bm{v}_2 = c_2 \begin{bmatrix} 2 \\ 3 - i\sqrt{7} \end{bmatrix}, \quad c_2 \neq 0
\end{equation}

\section{验证}

验证特征值和特征向量的关系$\bm{A}\bm{v} = \lambda\bm{v}$。

\subsection{验证$\lambda_1$和$\bm{v}_1$}

\begin{align}
\bm{A}\bm{v}_1 &= \begin{bmatrix} 1 & 1 \\ -4 & 4 \end{bmatrix} \begin{bmatrix} 2 \\ 3 + i\sqrt{7} \end{bmatrix} \\
&= \begin{bmatrix} 2 + 3 + i\sqrt{7} \\ -8 + 4(3 + i\sqrt{7}) \end{bmatrix} \\
&= \begin{bmatrix} 5 + i\sqrt{7} \\ -8 + 12 + 4i\sqrt{7} \end{bmatrix} \\
&= \begin{bmatrix} 5 + i\sqrt{7} \\ 4 + 4i\sqrt{7} \end{bmatrix}
\end{align}

\begin{align}
\lambda_1 \bm{v}_1 &= \frac{5 + i\sqrt{7}}{2} \begin{bmatrix} 2 \\ 3 + i\sqrt{7} \end{bmatrix} \\
&= \begin{bmatrix} 5 + i\sqrt{7} \\ \frac{(5 + i\sqrt{7})(3 + i\sqrt{7})}{2} \end{bmatrix} \\
&= \begin{bmatrix} 5 + i\sqrt{7} \\ \frac{15 + 5i\sqrt{7} + 3i\sqrt{7} + i^2 \cdot 7}{2} \end{bmatrix} \\
&= \begin{bmatrix} 5 + i\sqrt{7} \\ \frac{15 + 8i\sqrt{7} - 7}{2} \end{bmatrix} \\
&= \begin{bmatrix} 5 + i\sqrt{7} \\ \frac{8 + 8i\sqrt{7}}{2} \end{bmatrix} \\
&= \begin{bmatrix} 5 + i\sqrt{7} \\ 4 + 4i\sqrt{7} \end{bmatrix}
\end{align}

因此，$\bm{A}\bm{v}_1 = \lambda_1 \bm{v}_1$，验证正确！

\subsection{验证$\lambda_2$和$\bm{v}_2$}

\begin{align}
\bm{A}\bm{v}_2 &= \begin{bmatrix} 1 & 1 \\ -4 & 4 \end{bmatrix} \begin{bmatrix} 2 \\ 3 - i\sqrt{7} \end{bmatrix} \\
&= \begin{bmatrix} 2 + 3 - i\sqrt{7} \\ -8 + 4(3 - i\sqrt{7}) \end{bmatrix} \\
&= \begin{bmatrix} 5 - i\sqrt{7} \\ -8 + 12 - 4i\sqrt{7} \end{bmatrix} \\
&= \begin{bmatrix} 5 - i\sqrt{7} \\ 4 - 4i\sqrt{7} \end{bmatrix}
\end{align}

\begin{align}
\lambda_2 \bm{v}_2 &= \frac{5 - i\sqrt{7}}{2} \begin{bmatrix} 2 \\ 3 - i\sqrt{7} \end{bmatrix} \\
&= \begin{bmatrix} 5 - i\sqrt{7} \\ \frac{(5 - i\sqrt{7})(3 - i\sqrt{7})}{2} \end{bmatrix} \\
&= \begin{bmatrix} 5 - i\sqrt{7} \\ \frac{15 - 5i\sqrt{7} - 3i\sqrt{7} + i^2 \cdot 7}{2} \end{bmatrix} \\
&= \begin{bmatrix} 5 - i\sqrt{7} \\ \frac{15 - 8i\sqrt{7} - 7}{2} \end{bmatrix} \\
&= \begin{bmatrix} 5 - i\sqrt{7} \\ \frac{8 - 8i\sqrt{7}}{2} \end{bmatrix} \\
&= \begin{bmatrix} 5 - i\sqrt{7} \\ 4 - 4i\sqrt{7} \end{bmatrix}
\end{align}

因此，$\bm{A}\bm{v}_2 = \lambda_2 \bm{v}_2$，验证正确！

\section{总结}

矩阵$\bm{A} = \begin{bmatrix} 1 & 1 \\ -4 & 4 \end{bmatrix}$的特征值和特征向量为：

\textbf{特征值}：
\begin{align}
\lambda_1 &= \frac{5 + i\sqrt{7}}{2} \\
\lambda_2 &= \frac{5 - i\sqrt{7}}{2}
\end{align}

\textbf{对应的特征向量}：
\begin{align}
\bm{v}_1 &= \begin{bmatrix} 2 \\ 3 + i\sqrt{7} \end{bmatrix} \\
\bm{v}_2 &= \begin{bmatrix} 2 \\ 3 - i\sqrt{7} \end{bmatrix}
\end{align}

\textbf{注}：特征向量可以乘以任意非零标量，结果仍然是特征向量。

\end{document}

